\documentclass[a4paper,12pt]{article}
\usepackage[T2A]{fontenc}
\usepackage[utf8]{inputenc}
\usepackage[russian]{babel}
\usepackage{amsmath,amssymb,mathtools}
\usepackage{helvet}
\renewcommand{\familydefault}{\sfdefault}
\usepackage{icomma}
\usepackage[a4paper,margin=18mm,headheight=15pt]{geometry}
\usepackage{microtype}
\usepackage{tikz}
\usetikzlibrary{calc,arrows.meta,positioning,shapes.geometric}
\usepackage{tabularx,booktabs,longtable,array}
\usepackage{enumitem}
\usepackage{hyperref}
\usepackage{parskip}
\usepackage{fancyhdr}
\usepackage{setspace}
\usepackage{xcolor}
\usepackage{multicol}

\definecolor{accent}{HTML}{176B6B}
\definecolor{darkaccent}{HTML}{0D4444}
\definecolor{textgray}{HTML}{2D3338}
\definecolor{softgray}{HTML}{EEF2F2}
\definecolor{warn}{HTML}{8A4A32}

\hypersetup{colorlinks=true,linkcolor=darkaccent,urlcolor=darkaccent}
\setlength{\parindent}{0pt}
\setlength{\parskip}{0.55em}
\onehalfspacing
\pagestyle{fancy}
\fancyhf{}
\lhead{\small\color{textgray}Степени: свойства и преобразования}
\rhead{\small\color{textgray}Анна Баннова}
\cfoot{\small\color{textgray}\thepage}
\setlist[itemize]{leftmargin=1.4em,itemsep=0.25em,topsep=0.25em}
\setlist[enumerate]{leftmargin=1.7em,itemsep=0.45em,topsep=0.35em}
\newcommand{\term}[1]{\textbf{\color{darkaccent}#1}}
\newcommand{\note}[1]{\par\smallskip\noindent{\color{warn}\textbf{Проверка:}} #1\par\smallskip}
\newcommand{\formula}[1]{\[\boxed{\displaystyle #1}\]}

\tikzset{
  startstop/.style={ellipse,draw=accent,line width=0.9pt,minimum width=36mm,minimum height=10mm,align=center,fill=softgray},
  action/.style={rectangle,rounded corners=2pt,draw=accent,line width=0.9pt,text width=55mm,minimum height=12mm,align=center},
  decision/.style={diamond,aspect=2.25,draw=accent,line width=0.9pt,text width=33mm,align=center,inner sep=1.5pt},
  arrow/.style={-{Latex[length=2.6mm]},line width=0.9pt,draw=darkaccent},
  lab/.style={font=\small,fill=white,inner sep=1.5pt,text=textgray}
}

\begin{document}

% ---------- TITLE ----------
\thispagestyle{empty}
\begin{center}
\vspace*{24mm}
{\Large\color{textgray}Чек-лист}\par
\vspace{5mm}
{\Huge\bfseries\color{darkaccent}Степени}\par
\vspace{3mm}
{\LARGE\color{accent}свойства, границы применимости и перевод в степенной вид}\par
\vspace{14mm}
{\large Анна Баннова}\par
\vspace{2mm}
{\normalsize 2 октября 2026}\par
\vfill
{\large\bfseries Лёвин Артём Александрович}\par
{\normalsize эксклюзивно для Анны Банновой}\par
\vspace{22mm}
\begin{tikzpicture}[remember picture,overlay]
  \draw[accent,line width=1.1pt]
  ($(current page.south west)+(18mm,18mm)$)
  .. controls ($(current page.south west)+(65mm,33mm)$) and ($(current page.south east)+(-70mm,7mm)$) ..
  ($(current page.south east)+(-18mm,20mm)$);
\end{tikzpicture}
\end{center}
\clearpage

% ---------- CORE IDEAS ----------
\section*{1. Главная идея занятия}

В задачах со степенями сначала ищите структуру выражения. Формула работает только при выполнении своих условий. Одну и ту же формулу полезно узнавать \term{в обе стороны}: сворачивать выражение и, при необходимости, раскрывать его обратно.

\note{Плюс или минус между частями выражения часто означает, что стандартная формула для произведения или частного степеней здесь напрямую не применяется. Например, $(a+b)^n\ne a^n+b^n$ в общем случае.}

\subsection*{Базовые факты}
\begin{itemize}
  \item $a^1=a$.
  \item $a^0=1$ при $a\ne0$.
  \item $0^n=0$ при натуральном $n>0$; выражение $0^0$ в школьной алгебре не определяют.
\end{itemize}

\section*{2. Восемь формул, которые нужно узнавать}

\renewcommand{\arraystretch}{1.45}
\begin{tabularx}{\linewidth}{@{}>{\bfseries\color{darkaccent}}p{7mm} X X@{}}
\toprule
№ & Формула & Когда применять \\
\midrule
1 & $a^m\cdot a^n=a^{m+n}$ & одинаковое основание, между степенями умножение \\
2 & $\dfrac{a^m}{a^n}=a^{m-n}$, $a\ne0$ & одинаковое ненулевое основание, деление \\
3 & $a^n b^n=(ab)^n$ & одинаковая степень, между основаниями умножение \\
4 & $\dfrac{a^n}{b^n}=\left(\dfrac ab\right)^n$, $b\ne0$ & одинаковая степень в числителе и знаменателе \\
5 & $(a^m)^n=a^{mn}$ & степень возводится в степень \\
6 & $\dfrac1{a^n}=a^{-n}$, $a\ne0$ & нужно убрать дробь вида $1/a^n$ \\
7 & $\left(\dfrac ab\right)^n=\left(\dfrac ba\right)^{-n}$, $a,b\ne0$ & нужно поменять местами числитель и знаменатель \\
8 & $\sqrt[k]{a^m}=a^{m/k}$ & переход от корня к дробной степени; для единой школьной записи безопасно считать $a>0$ \\
\bottomrule
\end{tabularx}

\vspace{2mm}
\textbf{Формулы работают и справа налево.} Например,
\[
a^{m+n}=a^m\cdot a^n,\qquad (ab)^n=a^n b^n.
\]

\clearpage

% ---------- POWER FORM ----------
\section*{3. Принцип: переводим основание в степенной вид}

Когда получили задачу на степени, проверьте основание по трём пунктам.

\subsection*{1. Составное число}
\term{Простое число} имеет ровно два положительных делителя: $1$ и само число.\par
\term{Составное число} имеет больше двух положительных делителей. Число $1$ не является ни простым, ни составным.

Составное основание полезно разложить на простые множители:
\[
16=2^4,\qquad 81=3^4,\qquad 25=5^2,\qquad 49=7^2.
\]

\subsection*{2. Корень в основании}
Корень можно заменить дробной степенью:
\[
\sqrt[k]{a^m}=a^{m/k}.
\]
Примеры:
\[
\sqrt{5}=5^{1/2},\qquad \sqrt[3]{81}=\left(3^4\right)^{1/3}=3^{4/3}.
\]

\subsection*{3. Дробь в основании}
Дробь часто удобно заменить отрицательной степенью:
\[
\frac1{2^5}=2^{-5},\qquad \frac15=5^{-1},\qquad 0{,}5=\frac12=2^{-1}.
\]

\note{Этот алгоритм относится прежде всего к \textbf{основанию} степени. Например, в $3^{\sqrt5}$ корень находится в показателе; специально переписывать $\sqrt5$ как степень для этого алгоритма не требуется.}

\section*{4. Как понять, что преобразование законно}
{\small
\begin{itemize}
  \item Определите операцию между степенями: умножение, деление, сложение или вычитание.
  \item Проверьте, что совпадает именно то, что требует формула: основание или показатель.
  \item При делении проверьте ненулевой знаменатель; при корнях и дробных степенях — область допустимых значений.
  \item После шага проверьте: выражение стало проще и появилась ли следующая узнаваемая формула?
\end{itemize}
}

\clearpage

% ---------- FLOWCHART ----------
\thispagestyle{plain}
\begin{center}
{\Large\bfseries\color{darkaccent}Алгоритм: задача на вычисление со степенями}
\end{center}
\vspace{8mm}

\begin{center}
\begin{tikzpicture}[node distance=7mm and 15mm]
  \node[startstop] (s) {Получено выражение со степенями};
  \node[decision,below=of s] (comp) {Составное число в основании?};
  \node[action,right=of comp,text width=44mm] (factor) {Разложить основание на простые множители};
  \node[decision,below=11mm of comp] (root) {Корень в основании?};
  \node[action,right=of root,text width=44mm] (toroot) {Заменить корень дробной степенью};
  \node[decision,below=11mm of root] (frac) {Дробь в основании?};
  \node[action,right=of frac,text width=44mm] (tofrac) {При необходимости перейти к отрицательной степени};
  \node[action,below=12mm of frac,text width=72mm] (scan) {Сравнить выражение с формулами: одинаковые основания? одинаковые показатели? степень в степени?};
  \node[decision,below=of scan] (can) {Есть применимая формула?};
  \node[action,below=of can,text width=58mm] (apply) {Применить формулу и упростить};
  \node[action,right=of can,text width=44mm] (check) {Проверить условия и выбрать другой ход};
  \node[startstop,below=of apply] (finish) {Получено простое выражение / число};

  \draw[arrow] (s) -- (comp);
  \draw[arrow] (comp) -- node[lab,above]{да} (factor);
  \draw[arrow] (factor.south) |- (root.east);
  \draw[arrow] (comp) -- node[lab,left]{нет} (root);
  \draw[arrow] (root) -- node[lab,above]{да} (toroot);
  \draw[arrow] (toroot.south) |- (frac.east);
  \draw[arrow] (root) -- node[lab,left]{нет} (frac);
  \draw[arrow] (frac) -- node[lab,above]{да} (tofrac);
  \draw[arrow] (tofrac.south) |- (scan.east);
  \draw[arrow] (frac) -- node[lab,left]{нет} (scan);
  \draw[arrow] (scan) -- (can);
  \draw[arrow] (can) -- node[lab,left]{да} (apply);
  \draw[arrow] (can) -- node[lab,above]{нет} (check);
  \draw[arrow] (check.east) -- ++(9mm,0) |- (scan.east);
  \draw[arrow] (apply) -- (finish);
\end{tikzpicture}
\end{center}
\clearpage

% ---------- EXAMPLES ----------
\section*{5. Примеры по алгоритму}

\subsection*{Пример 1. Одинаковые основания}
\[
\frac{2^3\cdot2^4}{2^7}
=2^{3+4-7}
=2^0
=1.
\]
\textbf{Что сработало:} при умножении показатели сложили, при делении вычли.

\subsection*{Пример 2. Составное основание + степень в степени}
\[
5^{0{,}36}\cdot25^{0{,}32}
=5^{0{,}36}\cdot\left(5^2\right)^{0{,}32}
=5^{0{,}36}\cdot5^{0{,}64}
=5^1
=5.
\]
\textbf{Что сработало:} $25=5^2$, затем показатели при степени в степени перемножили, а при умножении степеней сложили.

\subsection*{Пример 3. Дробные показатели}
\[
7^{4/9}\cdot49^{5/18}
=7^{4/9}\cdot\left(7^2\right)^{5/18}
=7^{4/9}\cdot7^{5/9}
=7.
\]

\subsection*{Пример 4. Корень и дробь}
\[
\frac1{\sqrt3}
=\left(3^{1/2}\right)^{-1}
=3^{-1/2}.
\]
Здесь можно идти и в другом порядке: сначала $1/x=x^{-1}$, затем заменить корень дробной степенью. Результат совпадёт.

\section*{6. Частые ошибки}
\begin{enumerate}
  \item \textbf{Складывать показатели в $(a^m)^n$.} Здесь показатели \term{перемножаются}: $(a^m)^n=a^{mn}$.
  \item \textbf{Применять формулу через плюс:} $a^n+b^n\ne(a+b)^n$ в общем случае.
  \item \textbf{Забывать условие $a\ne0$:} оно необходимо для $a^0=1$ и формул с делением на $a^n$.
  \item \textbf{Ошибаться при факторизации:} например, $9=3^2$, а не $3^3$.
  \item \textbf{Менять дробь местами без смены знака показателя:} $\left(\frac ab\right)^n=\left(\frac ba\right)^{-n}$.
  \item \textbf{Преобразовывать всё подряд.} Меняйте запись только тогда, когда это открывает путь к следующей формуле.
\end{enumerate}

\clearpage

% ---------- TRAINING ----------
\section*{7. Тренировочный блок — Путь А}
\textbf{10 заданий.} Решайте по принципу: привести основания к удобному степенному виду, затем применить формулы.

\begin{enumerate}
  \item Вычислите: $2^4\cdot2^3$.
  \item Вычислите: $\dfrac{5^9}{5^6}$.
  \item Упростите: $\left(3^2\right)^4$.
  \item Укажите равенство, которое \textbf{не является тождеством}:\\
  \quad A) $x^4y^4=(xy)^4$;\qquad
  Б) $x^4+y^4=(x+y)^4$;\qquad
  В) $(x^4)^3=x^{12}$.
  \item Представьте в виде степени с основанием $2$: $\dfrac1{4^3}$.
  \item Представьте в виде степени с основанием $3$: $\sqrt[3]{81}$.
  \item Представьте в виде степени с основанием $5$: $\dfrac1{\sqrt5}$.
  \item Вычислите: $5^{0{,}4}\cdot25^{0{,}3}$.
  \item Вычислите: $3^{2/5}\cdot9^{3/10}$.
  \item Вычислите: $\dfrac{16^{3/4}}{8^{2/3}}$.
\end{enumerate}

\vfill
\begin{center}
\small\color{textgray}
Перед проверкой ответов убедитесь, что в каждой строке Вы можете назвать применённую формулу и её условие.
\end{center}
\clearpage

% ---------- ANSWERS ----------
\section*{8. Ответы}

\renewcommand{\arraystretch}{1.55}
\begin{tabularx}{\linewidth}{@{}>{\bfseries\color{darkaccent}}p{12mm} X@{}}
\toprule
№ & Ответ \\
\midrule
1 & $2^7=128$ \\
2 & $5^3=125$ \\
3 & $3^8=6561$ \\
4 & Б: $x^4+y^4=(x+y)^4$ — неверно в общем случае \\
5 & $2^{-6}$ \\
6 & $3^{4/3}$ \\
7 & $5^{-1/2}$ \\
8 & $5$ \\
9 & $3$ \\
10 & $2$ \\
\bottomrule
\end{tabularx}

\vfill
\begin{center}
\color{textgray}\small
Лёвин Артём Александрович эксклюзивно для Анны Банновой
\end{center}

\end{document}
